% ===================== CHAPTER 5 — SOLUTION AND CONTRIBUTION =====================
% Length: >= 5 pages. THE chapter graders weight most. Present each contribution in a
% section with: (i) problem, (ii) solution, (iii) results. Do NOT repeat earlier detail
% verbatim — earlier chapters should cross-reference here.
This chapter presents the principal contributions of the thesis. Each contribution
is described through the problem it addresses, the proposed solution, and the
results achieved.

\section{FEFO Lot-Allocation Mechanism}
% (i) Problem: perishable goods must leave by earliest expiry to reduce waste, while
%     keeping a per-lot audit trail. A plain quantity-only stock model cannot do this.
% (ii) Solution: an InventoryLot model + a FEFO query (ORDER BY expiryDate ASC NULLS
%     LAST) that greedily consumes the earliest-expiring lots on goods issue, recording
%     the batch/expiry on each issue line. (Reference: GoodsIssueServiceImpl.)
% (iii) Results: oldest stock exits first; full traceability; expiry alerts.
\thesisfig{flowchart_fefo}{FEFO allocation algorithm}{fig:fefo}
The first contribution is the lot-allocation mechanism for perishable goods.
For date-sensitive stock, inventory theory recommends issuing the
earliest-expiring units first so that shrinkage from expiry is
minimised~\cite{silver1998inventory, richards2021warehouse}. Commercial
distribution systems therefore provide batch and expiry tracking, for example the
lot and removal-strategy features of Odoo~\cite{odooFefo} and the batch
management of SAP Business One~\cite{sapBatch}. Many lighter-weight systems,
however, keep only an aggregate quantity for each product and warehouse. This
model is sufficient for simple stock checking, but it cannot answer which batch
should be picked first, which expiry date was delivered to a customer, or which
remaining lots are close to expiry. For a distributor that handles food,
medicine, cosmetics, or other date-sensitive goods, this limitation directly
creates waste and traceability risk.

The proposed solution extends the aggregate inventory model with an
\texttt{InventoryLot} entity. Each lot stores the product, warehouse, lot number,
expiry date, quantity received, quantity remaining, unit cost, and the source
goods receipt item. When a goods receipt is confirmed, the accepted quantity
increases aggregate inventory and also creates a corresponding inventory lot.
The implementation prevents duplicate lot creation for the same receipt item by
checking whether a source receipt item has already generated a lot. This keeps
goods receipt confirmation safe when the operation is retried after an
interrupted request.

During goods issue confirmation, the service retrieves all available lots for
the product and warehouse ordered by FEFO: lots with an expiry date come first,
the earliest expiry date is selected before later ones, and lots without expiry
dates are placed last. The service validates that the sum of remaining lot
quantity can cover the requested issue quantity. It then greedily consumes lots
in this order until the requested quantity is fully allocated. The first consumed
lot number and expiry date are copied to the goods issue item so that the
outbound document displays the batch information used for picking.

\begin{algorithm}[H]
\caption{FEFO allocation for a goods-issue item}
\KwIn{product, warehouse, issueQuantity}
\KwOut{updated lots and goods-issue line batch information}
lots $\leftarrow$ find available lots ordered by expiry date ascending\;
\If{sum(lot.quantityRemaining) $<$ issueQuantity}{
  reject confirmation with an insufficient-lot-stock error\;
}
remaining $\leftarrow$ issueQuantity\;
\ForEach{lot in lots}{
  \If{remaining = 0}{break}
  take $\leftarrow$ min(remaining, lot.quantityRemaining)\;
  lot.quantityRemaining $\leftarrow$ lot.quantityRemaining - take\;
  remaining $\leftarrow$ remaining - take\;
}
record the first consumed lot number and expiry date on the goods-issue item\;
\end{algorithm}

The achieved result is a practical FEFO process integrated with the order and
warehouse workflow. The system no longer relies on the warehouse employee to
manually choose a batch after reading several spreadsheets. Earlier-expiring
stock is consumed first, remaining lot quantities are visible for expiry
monitoring, and outbound documents keep the batch information needed for later
traceability.

% You may also typeset the algorithm with algorithm2e, e.g.:
% \begin{algorithm}[H]\caption{FEFO allocation}\dots\end{algorithm}

\section{Concurrency Control for Lot-Level Stock Deduction}
% (i) Problem: the FEFO lot deduction is a read-modify-write on InventoryLot, which has
%     no version/row lock. Two simultaneous goods-issue confirmations on the same scarce
%     product can both pass validation and both deduct -> lost update at lot level; the
%     per-lot ledger over-draws and diverges from the aggregate.
% (ii) Solution: acquire a pessimistic write lock (SELECT ... FOR UPDATE) on the single
%      aggregate Inventory row at the start of confirm(), before any lot is read, ordered
%      by ascending productId to prevent deadlock. (lockInventoryForUpdate.)
% (iii) Results: the conflicting confirmations are serialized at the database; the lot
%       ledger stays consistent with the aggregate.
The next contribution hardens the goods-issue path against a concurrency
hazard in the lot deduction described earlier in this chapter. The FEFO
allocation is a read-modify-write sequence: the service reads the remaining
quantity of each lot, computes how much to take, and writes the reduced
quantity back. The aggregate \texttt{Inventory} row is already protected, both
by an optimistic version field and by a pessimistic read in
\texttt{decreaseInventory} that issues a \texttt{SELECT~\dots~FOR~UPDATE}. The
per-lot \texttt{InventoryLot} entity, however, carries neither a version column
nor a row lock. When two warehouse staff confirm goods issues for the same
scarce product at almost the same time, both transactions can read the same lot
quantities, both can pass the validation that the remaining lot quantity covers
the requested amount, and both can deduct from the same lots. The second write
silently overwrites the first, which is the classic lost-update anomaly. The
result is a per-lot ledger that over-draws a lot or drives its remaining
quantity negative, so the sum of lot quantities no longer reconciles with the
aggregate on-hand quantity, defeating the traceability the lot model was built
to provide.

The proposed solution serializes the entire deduction of a product behind the
aggregate inventory row, which is the only record shared by every goods issue
for that product and warehouse. At the very start of the confirmation, before
any lot is read, the service acquires a pessimistic write lock on each
relevant \texttt{Inventory} row through a dedicated
\texttt{lockInventoryForUpdate} operation. This operation reuses the existing
\texttt{SELECT~\dots~FOR~UPDATE} query but does not modify the row; its only
purpose is to claim the lock. Because the lock is taken inside the
confirmation's transaction, the database holds it until that transaction
commits or rolls back~\cite{postgresql14}. A competing confirmation for the
same product therefore blocks on this single row, and when it finally proceeds
it re-reads the lots with the already-updated quantities rather than the stale
ones. In effect, the aggregate row acts as a gate that protects not only the
aggregate deduction but the unguarded lot read-modify-write as well. This is an
application of the pessimistic offline lock pattern, in which a shared record is
locked for the whole duration of a business transaction to prevent concurrent
modification~\cite{fowler2002poeaa, bauer2015hibernate}.

A goods issue may contain several products, so a single confirmation can lock
several inventory rows. To avoid a deadlock in which one transaction holds the
lock on product A while waiting for product B and a second transaction holds B
while waiting for A, the service collects the distinct product identifiers and
acquires the locks in ascending identifier order. Because every confirmation
takes the locks in the same global order, a cycle of mutual waiting cannot
form, which is the standard ordered-locking discipline for deadlock
prevention~\cite{fowler2002poeaa}.

The result is that two conflicting goods-issue confirmations for the same
product no longer interleave; the second is queued at the database until the
first completes, after which it observes a consistent state. The lot ledger
stays reconciled with the aggregate on-hand quantity even under concurrent
issue, which closes the lost-update gap while keeping the database schema
unchanged, since the fix reuses the existing aggregate row rather than adding a
version column to every lot. The accepted trade-off is reduced parallelism for
issues that touch the same product; this is acceptable because such operations
are inherently conflicting and must be ordered for correctness. This behaviour
is guarded by a deterministic two-thread integration test
(\texttt{ConcurrencyLotTest}) in which both confirmations are released
simultaneously through a shared barrier: with the lock in place exactly one
issue succeeds and the lot ledger stays consistent with the aggregate on-hand
quantity, whereas removing the lock reintroduces the lost update.

\section{Inventory Reservation on Order Approval}
% (i) Problem: overselling between approval and shipment.
% (ii) Solution: reserve stock at SO approval (quantityReserved up, quantityAvailable
%      down); deduct only at goods-issue confirmation; release on cancel.
% (iii) Results: no oversell; clear available-vs-reserved visibility.
The second contribution is inventory reservation at the moment a sales order is
approved. In a manual process, two sales staff members may check the same
physical stock and both promise it to different customers. Even if the stock is
physically sufficient at the first check, the company can oversell before the
warehouse confirms goods issue. This problem appears in the interval between
sales approval and physical shipment, and it corresponds to the
available-to-promise concept in supply chain planning, where committed stock must
be separated from free stock before an order is confirmed~\cite{chopra2019supply}.

The system addresses this issue by separating on-hand, reserved, and available
quantities. When a sales order is approved, the service checks available
quantity for each order line in the selected warehouse. If every line can be
fulfilled, it increases the reserved quantity and reduces the available quantity.
The physical on-hand quantity is not reduced at this stage because the goods are
still in the warehouse. When the goods issue is later confirmed, the system
deducts on-hand quantity and releases the corresponding reservation. If an
approved order is cancelled before shipment, only the remaining reserved
quantity is released.

This design gives different business users the right view of stock. Sales staff
see available quantity, which prevents them from selling stock that is already
committed. Warehouse staff still see the physical on-hand quantity for picking
and reconciliation. Managers can distinguish between real stock shortage and
stock that is present but reserved for approved orders. As a result, the system
reduces overselling risk without prematurely deducting stock before goods leave
the warehouse.

\section{Failed-Delivery Recovery}
% (i) Problem: a failed delivery must cleanly reverse stock, invoice and order state.
% (ii) Solution: restockFromFailedDelivery() restores inventory + lots, cancels the
%      invoice (warns if already paid), resets delivered qty, reverts SO status.
% (iii) Results: consistent, idempotent recovery.
The third contribution is a centralized recovery mechanism for failed delivery.
In the order-to-cash workflow, a confirmed goods issue has already deducted
stock, updated delivered quantities, and generated an invoice. If the delivery
later fails, a manual correction must reverse several records in the right order.
If only one record is corrected, the system becomes inconsistent: inventory may
show fewer goods than the warehouse has, the order may appear delivered, or the
invoice may remain active although the customer did not receive the goods.

The system implements this recovery in \texttt{restockFromFailedDelivery}. The
method finds the goods issue linked to the failed trip, adds the issued quantity
back to aggregate inventory, restores the quantity to the corresponding lot,
decreases the delivered quantity of the sales order item, cancels the invoice
when it is still cancellable and unpaid, and recalculates the sales order
delivery status. If the invoice already has payment, the system does not
silently cancel it; instead it records a warning so that the accountant can
perform refund or adjustment work manually. This is important because financial
documents should not be reversed automatically after money has been collected.

The result is that delivery failure is treated as a business event rather than a
set of unrelated manual edits. The recovery path protects the consistency among
inventory, lots, sales orders, and invoices. In the current implementation, the
method is intended to be called from a controlled delivery-status transition; a
future improvement is to add an explicit recovery marker on the goods issue or
delivery order to make repeated calls harmless at the database level.

\section{Role-Based Workflow Enforcement}
% (i) Problem: separation of duties (creator != approver) across 9 roles.
% (ii) Solution: RBAC at endpoint + method + data level (e.g. shippers see only their
%      own trips via isAssignedTo()).
% (iii) Results: enforced approval chains and least-privilege access.
The fourth contribution is role-based workflow enforcement. The distribution
process involves several responsibilities: staff create documents, managers or
accountants approve them, warehouse staff confirm physical stock movements,
delivery administrators assign trips, and shippers update only their assigned
deliveries. If all users could access all operations, the system would lose
separation of duties and approval would become only a visual status rather than
a real control. This requirement is an instance of the role-based access control
model, in which permissions are assigned to roles rather than to individual
users, and which supports separation of duties and least
privilege~\cite{sandhu2000nist}.

The system defines nine roles and applies access control at three levels. At the
endpoint level, Spring Security restricts URL patterns and HTTP methods, for
example allowing only purchase managers, accountants, or administrators to
approve purchase orders. At the method level, \texttt{@PreAuthorize} annotations on individual
controller endpoints restrict sensitive operations, such as approving and
rejecting orders, recording payments, and managing user accounts, to the roles
permitted to perform them. At the data level, services enforce contextual
constraints, such as preventing warehouse staff from seeing unapproved orders
and limiting shippers to assigned trips. Table~\ref{tab:rbac-matrix} summarises
the main permission groups.

\begin{table}[H]\centering
\caption{Simplified RBAC permission matrix}
\label{tab:rbac-matrix}
\small
\resizebox{\linewidth}{!}{%
\begin{tabular}{|p{3.6cm}|c|c|c|c|c|c|}
\hline
\textbf{Role} & \textbf{PO} & \textbf{SO} & \textbf{GR/GI} & \textbf{Delivery} & \textbf{Invoice} & \textbf{Users} \\ \hline
Administrator & Full & Full & Full & Full & Full & Full \\ \hline
Purchase staff & Create/Edit & View & View GR & No & No & No \\ \hline
Purchase manager & Approve & View & View GR & No & No & No \\ \hline
Sales staff & View & Create/Edit & View GI & No & View & No \\ \hline
Sales manager & View & Approve & View GI & No & Manage & No \\ \hline
Warehouse staff & Approved only & Approved only & Confirm & No & No & No \\ \hline
Delivery administrator & No & View & View & Manage & No & No \\ \hline
Shipper & No & No & No & Assigned trips & No & No \\ \hline
Accountant & Approve & Approve & View & No & Manage & No \\ \hline
\end{tabular}%
}
\end{table}

This design makes the workflow enforceable by the application instead of relying
only on organizational discipline. It also improves usability because each user
sees a smaller set of relevant actions. The result is a least-privilege
prototype that supports approval chains, warehouse control, delivery assignment,
and financial visibility while keeping administrative functions restricted.

\section{Eliminating N+1 Queries in Delivery Listing}
% (i) Problem: object-relational N+1 — one extra query per row + lazy loads made the
%     available-delivery-order listing scale linearly with data and become slow.
% (ii) Solution: replace per-row queries with three bulk queries (orders mapped by code,
%     planned order ids as a set, confirmed goods issues with fetch-joined associations).
% (iii) Results: endpoint latency 2380 ms -> 47 ms on the same dataset (~50x).
The fifth contribution is a performance optimisation of the screen that lists the
delivery orders available for planning. This endpoint is opened frequently by the
delivery administrator, so its responsiveness affects daily operation directly.
The original implementation iterated over every confirmed goods issue and, for
each one, executed a separate query to look up the matching delivery order by
code and another query to check whether that order already belonged to a delivery
plan. Because the sales order, customer, and delivery address of each goods issue
were mapped lazily, accessing them inside the loop triggered further queries. This
is the well-known object-relational \emph{N+1 select} problem documented in the
object-relational mapping literature~\cite{bauer2015hibernate}: the number of
database round trips grows linearly with the number of rows, so the response time
degrades as the data set grows. On the development data set the listing endpoint
took approximately 2380~ms.

The optimisation replaces the per-row queries with three bulk queries executed
once before the loop. The first loads every delivery order and indexes it by code
in an in-memory map, so the matching order is resolved by a hash lookup instead of
a database query. The second loads the identifiers of all delivery orders that
already belong to a plan as a single projection and stores them in a set, so the
``already planned'' test becomes a set membership check. The third loads the
confirmed goods issues with their sales order, customer, and delivery address
eagerly through a \texttt{JOIN FETCH} query, which removes the lazy-loading round
trips inside the loop. The loop then only creates a delivery order in the rare
case that one has not yet been materialised for a newly confirmed goods issue.

The result is a substantial latency reduction: on the same data set and
environment the endpoint response time fell from approximately 2380~ms to
47~ms, an improvement of roughly fifty times, while the returned data is
unchanged. The number of database round trips no longer depends on the number of
goods issues, so the listing remains responsive as the volume of confirmed
shipments grows. This contribution demonstrates that the system was not only
built to be functionally correct but was also profiled and tuned at the data
access layer, which is a common source of performance problems in applications
built on an object-relational mapping framework such as Spring Data JPA and
Hibernate~\cite{springboot314}.
